% SPDX-FileCopyrightText: Copyright (C) Nile Jocson <atraphaxia@gmail.com>
% SPDX-License-Identifier: MPL-2.0



\documentclass{article}



\usepackage{adjustbox}
\usepackage{amsmath}
\usepackage{circuitikz}
\usepackage{empheq}
\usepackage[a4paper, margin=1in]{geometry}
\usepackage{siunitx}



\ctikzset{american}



\title{
	EEE 123 \\
	Exercise 2
}

\author{
	Nile Jocson \\
	\textless{}atraphaxia@gmail.com\textgreater{}
}




\begin{document}
	\maketitle{}
		\pagebreak{}

	\tableofcontents{}
		\pagebreak{}

	\section{}
		Show your solutions on the space provided. You may use the back of the page as
		continuation of your solution. Please write legibly. \textbf{Box your final
		answers.} Use at least 2 decimal places if the answer is fractional.



		\subsection{}
			Use a combination of source transformation and circuit equivalence to find $V$
			and $I$ in the circuit below.
			\begin{center}\begin{circuitikz}
				\draw(0, 0)
					to[short, -*] ++(-3, 0)
					to[short, -*] ++(-3, 0)
					to[short] ++(-3, 0)
					to[V, invert, a = \qty{18}{\volt}] ++(0, 3)
					to[R, l = \qty{3}{\ohm}, -*] ++(3, 0)
					to[short, -*] ++(3, 0)
					to[R, l = \qty{4}{\ohm}, i = $I$] ++(3, 0)
					to[R, l = \qty{8}{\ohm}, v = $V$] ++(0, -3);
				\draw(-3, 0)
					to[I, l = \qty{4}{\ampere}] ++(0, 3);
				\draw(-6, 0)
					to[R, a = \qty{6}{\ohm}] ++(0, 3);
			\end{circuitikz}\end{center}

			Transform the \qty{18}{\volt} source.
				\[I_{\text{eq}} = \frac{18}{3} = 6\]

			\begin{center}\begin{circuitikz}
				\draw(0, 0)
					to[short, -*] ++(-3, 0)
					to[short, -*] ++(-3, 0)
					to[short, -*] ++(-3, 0)
					to[short] ++(-3, 0)
					to[I, a = \qty{6}{\ampere}] ++(0, 3)
					to[short, -*] ++(3, 0)
					to[short, -*] ++(3, 0)
					to[short, -*] ++(3, 0)
					to[R, l = \qty{4}{\ohm}, i = $I$] ++(3, 0)
					to[R, l = \qty{8}{\ohm}, v = $V$] ++(0, -3);
				\draw(-3, 0)
					to[I, l = \qty{4}{\ampere}] ++(0, 3);
				\draw(-6, 0)
					to[R, a = \qty{6}{\ohm}] ++(0, 3);
				\draw(-9, 0)
					to[R, a = \qty{3}{\ohm}] ++(0, 3);
			\end{circuitikz}\end{center}

			Combine parallel current sources and resistors.
				\[I_{\text{eq}} = 6 + 4 = 10\]
				\[R_{\text{eq}} = {\left(\frac{1}{6} + \frac{1}{4}\right)}^{-1} = 2\]

			\begin{center}\begin{circuitikz}
				\draw(0, 0)
					to[short, -*] ++(-3, 0)
					to[short] ++(-3, 0)
					to[I, a = \qty{10}{\ampere}] ++(0, 3)
					to[short, -*] ++(3, 0)
					to[R, l = \qty{4}{\ohm}, i = $I$] ++(3, 0)
					to[R, l = \qty{8}{\ohm}, v = $V$] ++(0, -3);
				\draw(-3, 0)
					to[R, l = \qty{2}{\ohm}] ++(0, 3);
				\draw[dashed, color = red] (-2.75, 3.75)
					-- (0.75, 3.75)
					-- (0.75, -0.25)
					-- (-0.75, -0.25)
					-- (-0.75, 2.25)
					-- (-2.75, 2.25)
					-- cycle;
				\draw[color = red] (0.75, 3.75) node[label = $\qty{12}{\ohm}$] {};
			\end{circuitikz}\end{center}

			Solve for $I$:
				\[I = 10\left(\frac{\frac{1}{12}}{\frac{1}{12} + \frac{1}{2}}\right) = \frac{10}{7}\]

			Solve for $V$:
				\[V = 8\left(\frac{10}{7}\right) = \frac{80}{7}\]

			Final answer:
				\begin{empheq}[box=\fbox]{align*}
					V &= \qty{11.43}{\volt} \\
					I &= \qty{1.43}{\ampere}
				\end{empheq}

			\pagebreak{}
\end{document}
